\section{Método da Bissecção}
\label{sec:metodo_bisseccao}

O Método da Bissecção é uma abordagem numérica clássica e incondicionalmente convergente para a determinação de raízes reais de equações não lineares. Fundamentado no Teorema de Bolzano, o método opera mediante a divisão sucessiva de um intervalo fechado $[a, b]$, no qual a função contínua $f(x)$ apresenta sinais opostos em seus extremos ($f(a) \cdot f(b) < 0$). Embora possua uma taxa de convergência linear — exigindo um número razoável de iterações —, sua robustez o torna ideal para isolamento inicial de raízes ou resolução de problemas onde a estabilidade é prioritária em relação à velocidade computacional.

\subsection{Estratégia de Implementação}
\label{subsec:estrategia_implementacao_bisseccao}

A concepção computacional do algoritmo foi estruturada na linguagem Python, priorizando a estabilidade numérica e a modularidade do código. A implementação principal encontra-se encapsulada no módulo \texttt{src/metodos/zeros\_de\_funcao.py}.

\begin{lstlisting}[language=Python, caption={Código de implementação do Método da Bissecção}, label={lst:bisseccao_codigo}]
def bisseccao(funcao, a, b, tol=1e-6, max_iter=100):
    historico = [] 
    
    for i in range(max_iter):
        c = (a + b)/2
        fc = funcao(c)
        denominador = a + b

        if denominador != 0:
            er = abs((a - b) / denominador)
        else:
            er = abs(b - a) 
        
        historico.append({
            'iter': i, 'a': a, 'b': b, 'c': c, 
            'f(c)': fc, 'erro': er
        })
        
        if fc == 0 or er < tol:
            return c, i, historico
            
        if utils.teorema_bolzano(funcao, a, c):
            b = c
        else:
            a = c
            
    return c, max_iter, historico
\end{lstlisting}

Para a formulação matemática do erro interativo e visando superar limitações do ambiente de programação, estabeleceu-se uma métrica adaptativa. Originalmente, o erro relativo ponderado pela amplitude do intervalo é definido como $\epsilon_{rel} = \left| \frac{a_k - b_k}{a_k + b_k} \right|$. Contudo, constatou-se que intervalos simétricos em relação à origem ($a_k = -b_k$) induziriam o interpretador Python a lançar uma exceção de divisão por zero (\texttt{ZeroDivisionError}). Para mitigar este risco, implementou-se um fluxo condicional: caso o denominador seja nulo, a métrica transita automaticamente para o erro absoluto $\epsilon_{abs} = |b_k - a_k|$. 

Os critérios de parada foram definidos para garantir o término do laço em tempo finito, compreendendo:
\begin{enumerate}
    \item \textbf{Raiz exata:} Interrupção imediata caso a avaliação funcional no ponto médio resulte exatamente em zero ($f(c_k) = 0$), limitando processamento desnecessário por flutuação de ponto flutuante.
    \item \textbf{Tolerância (\texttt{tol}):} O método encerra com sucesso caso a métrica de erro calculada seja estritamente menor que a tolerância estipulada.
    \item \textbf{Esforço computacional (\texttt{max\_iter}):} Uma trava de segurança máxima de iterações para abortar o processo caso a convergência não seja alcançada no limite imposto.
\end{enumerate}

\subsection{Estrutura dos Arquivos de Entrada/Saída}
\label{subsec:estrutura_arquivos_bisseccao}

Visando aderência aos princípios de responsabilidade única e reprodutibilidade científica, a persistência de dados foi dissociada da lógica matemática. Os diretórios foram organizados hierarquicamente:

\begin{itemize}
    \item \textbf{Dados de Entrada (\texttt{dados/entradas/}):} Estruturados em arquivos de texto simples (\texttt{.txt}). Foram escolhidos por permitirem fácil edição manual dos parâmetros de controle (limites $a$ e $b$, tolerância e limite de iterações) sem necessidade de reconfiguração no código-fonte.
    \item \textbf{Dados de Saída (\texttt{dados/saida/<questao>/}):} Os resultados iterativos foram consolidados em arquivos de valores separados por vírgula (\texttt{.csv}). A escolha do formato CSV justifica-se por sua interoperabilidade universal, permitindo a ingestão direta por bibliotecas de análise de dados (como Pandas) e por pacotes de renderização gráfica em LaTeX (como TikZ/pgfplots), viabilizando análises quantitativas automatizadas e plots de convergência nativos no documento final.
\end{itemize}

\subsubsection{Parâmetros de Entrada para a Bissecção}
\label{subsubsec:parametros_entrada_bisseccao}

Os arquivos de entrada definem os hiperparâmetros críticos para cada execução do método. A Tabela \ref{tab:parametros_bisseccao} sumariza os parâmetros utilizados nas questões propostas:

\begin{table}[htpb]
\centering
\caption{Parâmetros de configuração para o Método da Bissecção.}
\label{tab:parametros_bisseccao}
\small
\begin{tabular}{@{}lccccc@{}}
\toprule
\textbf{Questão} & \textbf{Função} & \textbf{Intervalo [a, b]} & \textbf{Tolerância} & \textbf{Máx Iter} \\ \midrule
3.3 & $f_1(T)$ & [1.0, 2.0] & $1 \times 10^{-5}$ & 100 \\
3.3 & $f_2(T)$ & [5.0, 9.0] & $1 \times 10^{-5}$ & 100 \\
3.6 & $f(x)$ & [30.0, 60.0] & $1 \times 10^{-4}$ & 100 \\
\bottomrule
\end{tabular}
\end{table}

Observa-se que a tolerância de $10^{-5}$ na Questão 3.3 exigirá um número maior de iterações comparativamente à tolerância mais relaxada de $10^{-4}$ na Questão 3.6. O limite máximo de 100 iterações constitui uma salvaguarda conservadora, garantindo que o algoritmo não incorra em consumo excessivo de recursos computacionais.

\subsection{Problemas Teste}
\label{subsec:problemas_testes_bisseccao}

A verificação e validação da implementação computacional foram conduzidas por meio de problemas aplicados extraídos da literatura de referência de Cálculo Numérico (Franco, Neide). O objetivo desta seção é atestar o funcionamento, a eficiência e a validade matemática das soluções encontradas.

\subsubsection{Problema Teste 1: Tempo de Subida de um Amplificador}
\label{subsubsec:problema_teste_1_bisseccao}

O primeiro problema propõe a determinação do tempo de subida de um amplificador eletrônico, compreendido como o intervalo temporal necessário para que a resposta $g(T)$ avance de $10\%$ a $90\%$ de seu valor assintótico, sendo a função normalizada expressa por:
\begin{equation}
    g(T) = 1 - \left(1 + T + \frac{T^2}{2}\right)e^{-T}
\end{equation}

O cálculo do tempo de subida, $\Delta T = T_{0.9} - T_{0.1}$, exige a determinação das raízes das funções auxiliares:
\begin{align}
    f_1(T) &= 0.9 - \left(1 + T + \frac{T^2}{2}\right)e^{-T} = 0 \label{eq:f1_bisseccao} \\
    f_2(T) &= 0.1 - \left(1 + T + \frac{T^2}{2}\right)e^{-T} = 0 \label{eq:f2_bisseccao}
\end{align}

% \textbf{Aquivo de entrada da solução:} O arquivo \texttt{questao\_3\_3\_bisseccao.txt} contém os parâmetros de controle para a execução do método, incluindo os intervalos iniciais $[1.0, 2.0]$ para $f_1(T)$ e $[0.0, 1.0]$ para $f_2(T)$, a tolerância de $10^{-5}$ e o limite máximo de iterações.


O algoritmo processou $f_1(T)$ com intervalo inicial $[1.0, 2.0]$ e tolerância $10^{-5}$. A Tabela \ref{tab:bisseccao_f1} exibe a progressão numérica.

A convergência foi atingida na 16ª iteração, com o erro caindo abaixo da tolerância ($\epsilon_{16} \approx 6.92 \times 10^{-6}$), retornando $T_{0.1} \approx 1.102058$. A Figura \ref{fig:grafico_bisseccao_csv} ilustra este amortecimento da magnitude funcional, que demonstra a redução progressiva assintótica característica da contração de intervalo.

\begin{figure}[H]
    \centering
    \begin{tikzpicture}
        \begin{axis}[
            width=14cm, height=8cm,
            title={\bfseries Evolução do Valor do Resíduo \(f(c)\) no Método da Bissecção},
            xlabel={Iteração (\(k\))},
            ylabel={Valor de \(f_1(c_k)\)},
            grid=both,
            grid style={line width=.1pt, draw=gray!20},
            major grid style={line width=.2pt,draw=gray!50},
            legend pos=north east,
            colormap name=viridis,
            tick label style={font=\footnotesize},
            label style={font=\small},
        ]
        \addplot[
            color=magenta, mark=*, mark options={fill=white, scale=1.2}, thick 
        ] table [col sep=comma, x=iter, y={f(c)}] {../dados/saida/questao_3_3/questao_3_3_hist_f1_bisseccao.csv};
        \addlegendentry{Resíduo \(f_1(c)\)}
        \end{axis}
    \end{tikzpicture}
    \caption{Convergência temporal da magnitude do erro para $T_{0.1}$.}
    \label{fig:grafico_bisseccao_csv}
\end{figure}

\textbf{Validação da Solução:} Para testar se o valor numérico é de fato a solução, reintroduz-se o ponto obtido $T_{0.1} = 1.102058$ na equação original:
\begin{equation}
    f_1(1.102058) = 0.9 - \left(1 + 1.102058 + \frac{(1.102058)^2}{2}\right)e^{-1.102058} \approx -0.00000140
\end{equation}

\begin{table}[htpb]
\centering
\caption{Histórico iterativo do Método da Bissecção para a busca de $T_{0.1}$ na função $f_1(T)$.}
\label{tab:bisseccao_f1}
\small
\begin{tabular}{@{}cccccc@{}}
\toprule
\textbf{Iteração} & $\mathbf{a_k}$ & $\mathbf{b_k}$ & $\mathbf{c_k}$ & $\mathbf{f_1(c_k)}$ & \textbf{Erro} ($\epsilon_k$) \\ \midrule
0  & 1.000000 & 2.000000 & 1.500000 & \phantom{-}0.09115317 & 0.33333333 \\
1  & 1.000000 & 1.500000 & 1.250000 & \phantom{-}0.03153233 & 0.20000000 \\
2  & 1.000000 & 1.250000 & 1.125000 & \phantom{-}0.00466937 & 0.11111111 \\
3  & 1.000000 & 1.125000 & 1.062500 & -0.00785071           & 0.05882353 \\
$\vdots$ & $\vdots$ & $\vdots$ & $\vdots$ & $\vdots$          & $\vdots$   \\
14 & 1.102051 & 1.102112 & 1.102081 & \phantom{-}0.00000322 & 0.00002769 \\
15 & 1.102051 & 1.102081 & 1.102066 & \phantom{-}0.00000014 & 0.00001385 \\
16 & 1.102051 & 1.102066 & 1.102058 & -0.00000140           & 0.00000692 \\ \bottomrule
\end{tabular}
\end{table}

Observa-se que o resíduo apresenta magnitude estritamente compatível com a precisão exigida, confirmando tratar-se da raiz buscada. Procedimento análogo para $f_2(T)$ rendeu $T_{0.9} \approx 5.322320$. Consequentemente, o tempo de subida validado pelo método é $\Delta T = 5.322320 - 1.102058 = 4.220262$. Os resultados demonstram que a implementação comporta-se com a eficiência teórica esperada $\mathcal{O}(\log(1/\epsilon))$.

\subsubsection{Problema Teste 2: Ângulo de Lançamento Balístico de um Míssil}
\label{subsubsec:problema_teste_2_bisseccao}

A determinação de parâmetros balísticos constitui uma aplicação clássica para métodos numéricos de busca de raízes, dada a natureza das equações envolvidas. O Problema 3.6 demanda o cálculo do ângulo de inclinação ($\alpha$) para o lançamento de um míssil, de modo a atingir um alvo situado a um deslocamento angular $\theta$ (medido a partir do centro da Terra). A cinemática do problema é modelada pela equação:
\begin{equation}
    \operatorname{tg}\left(\frac{\theta}{2}\right) = \frac{\sin \alpha \cos \alpha}{\frac{g R}{v^2} - \cos^2 \alpha}
    \label{eq:missil_original}
\end{equation}

Para submeter a Equação \ref{eq:missil_original} ao Método da Bissecção é necessário manipulá-la para a forma padrão de determinação de zeros, $f(\alpha) = 0$. A fim de mitigar potenciais instabilidades numéricas oriundas de divisões por termos próximos a zero, optou-se pela linearização do denominador através de multiplicação cruzada. Considerando os dados fornecidos, tem-se o deslocamento do alvo $\theta = 80^\circ$ e a relação adimensional de energia $\frac{v^2}{gR} = 1.25$, o que implica diretamente que $\frac{gR}{v^2} = 0.8$. A função objetivo, avaliada estritamente em radianos para compatibilidade com as bibliotecas matemáticas, torna-se:
\begin{equation}
    f(\alpha) = \sin(\alpha)\cos(\alpha) - \operatorname{tg}\left(40^\circ\right)\left(0.8 - \cos^2(\alpha)\right) = 0
    \label{eq:missil_funcao_bisseccao}
\end{equation}

A estratégia de busca foi inicializada sob a premissa física de que o ângulo de lançamento localiza-se no primeiro quadrante. O intervalo submetido foi $[30^\circ, 60^\circ]$, cujos equivalentes em radianos correspondem a aproximadamente $[0.5236, 1.0472]$. O limite de tolerância estipulado para o erro relativo da amplitude do intervalo foi de $\epsilon < 10^{-4}$.

A Tabela \ref{tab:bisseccao_missil} apresenta o histórico da progressão geométrica do hipervolume de busca ao longo das iterações computadas.

\begin{table}[htpb]
\centering
\caption{Histórico iterativo do Método da Bissecção para a função do ângulo de lançamento balístico $f(\alpha)$.}
\label{tab:bisseccao_missil}
\small
\begin{tabular}{@{}cccccc@{}}
\toprule
\textbf{Iteração ($k$)} & $\mathbf{a_k}$ & $\mathbf{b_k}$ & $\mathbf{c_k}$ & $\mathbf{f(c_k)}$ & \textbf{Erro} ($\epsilon_k$) \\ \midrule
0  & 0.523598 & 1.047197 & 0.785398 & \phantom{-}0.24827011 & 0.33333333 \\
1  & 0.785398 & 1.047197 & 0.916297 & \phantom{-}0.12264554 & 0.14285714 \\
2  & 0.916297 & 1.047197 & 0.981747 & \phantom{-}0.04965511 & 0.06666666 \\
3  & 0.981747 & 1.047197 & 1.014472 & \phantom{-}0.01114434 & 0.03225806 \\
4  & 1.014472 & 1.047197 & 1.030835 & -0.00854367           & 0.01587301 \\
5  & 1.014472 & 1.030835 & 1.022653 & \phantom{-}0.00133421 & 0.00800000 \\
6  & 1.022653 & 1.030835 & 1.026744 & -0.00359642           & 0.00398406 \\
$\vdots$ & $\vdots$ & $\vdots$ & $\vdots$ & $\vdots$          & $\vdots$   \\
10 & 1.023676 & 1.024187 & 1.023932 & -0.00020480           & 0.00024968 \\
11 & 1.023676 & 1.023932 & 1.023804 & -0.00005083           & 0.00012485 \\
12 & 1.023676 & 1.023804 & 1.023740 & \phantom{-}0.00002614 & 0.00006243 \\ \bottomrule
\end{tabular}
\end{table}
\vspace{1cm}

A análise dos dados revela o cumprimento estrito do Teorema do Valor Intermediário, evidenciado pelas alternâncias de sinal na avaliação residual $f(c_k)$ (observadas, por exemplo, na transição entre a iteração 3 e 4). O algoritmo operou de maneira estável, sem deflagrar a condição de exceção arquitetada para divisões por zero, uma vez que o domínio da busca permaneceu estritamente positivo (no primeiro quadrante trigonométrico).

A convergência foi matematicamente estabelecida na 12ª iteração ($k=12$). Neste estágio, a métrica de erro relativo $\epsilon_{12} \approx 6.24 \times 10^{-5}$ tornou-se inferior à barreira de tolerância imposta. O método extraiu a raiz aproximada $c_{12} \approx 1.023740$ radianos. Este valor numérico converte-se em um ângulo de inclinação de lançamento $\alpha \approx 58.65^\circ$.

\textbf{Validação e Eficiência da Solução:}
Para atestar a validade da aproximação numérica, submete-se o valor $c_{12} = 1.023740$ à Equação \ref{eq:missil_funcao_bisseccao}. O resíduo operacionalizado, $f(c_{12}) \approx 2.61 \times 10^{-5}$, apresenta uma magnitude da mesma ordem do erro de truncamento exigido. Isso confirma que a raiz satisfaz o modelo balístico proposto, atestando a robustez da implementação do Método da Bissecção frente a problemas contendo composições de funções trigonométricas. 

\begin{figure}[H]
    \centering
    \begin{tikzpicture}
        \begin{axis}[
            width=14cm, height=7cm,
            title={\bfseries Amortecimento Residual para o Ângulo de Lançamento},
            xlabel={Iteração (\(k\))},
            ylabel={Magnitude do Erro \(\epsilon_k\) e Resíduo},
            ymode=log, % Escala logarítmica é ideal para demonstrar convergência linear
            grid=both,
            grid style={line width=.1pt, draw=gray!20},
            major grid style={line width=.2pt,draw=gray!50},
            legend pos=north east,
            tick label style={font=\footnotesize},
            label style={font=\small},
        ]
        
        % Plot do erro 
        \addplot[
            color=blue, mark=square, thick 
        ] table [col sep=comma, x=iter, y=erro] {../dados/saida/questao_3_6/questao_3_6_hist_f1_bisseccao.csv};
        \addlegendentry{Erro Relativo (\(\epsilon_k\))}

        % Plot do valor absoluto da função
        \addplot[
            color=magenta, mark=*, mark options={fill=white}, thick 
        ] table [col sep=comma, x=iter, y expr=abs(\thisrow{f(c)})] {../dados/saida/questao_3_6/questao_3_6_hist_f1_bisseccao.csv};
        \addlegendentry{\(|f(c_k)|\)}

        \end{axis}
    \end{tikzpicture}
    \caption{Convergência linear na busca de $\alpha$ (Problema 3.6).}
    \label{fig:convergencia_bisseccao_missil}
\end{figure}

A Figura \ref{fig:convergencia_bisseccao_missil} corrobora de forma visual a análise, demonstrando, através da escala semi-logarítmica, que o erro de fato diminui à taxa constante previsível pela relação de recorrência $|x^* - c_k| \leq \frac{b-a}{2^k}$. A eficiência do método para este problema em específico provou-se altamente satisfatória, isolando a raiz sem exigir recursos adicionais para mitigar pontos de inflexão que poderiam prejudicar métodos abertos baseados em derivadas.


% Para inclusão futura: Descrever o problema extraído do livro de Cálculo Numérico de Neide Franco.
% - Inserir a equação avaliada f(x) = 0.
% - Inserir o intervalo inicial e tolerância requerida.
% - Importar a Tabela CSV via \pgfplotstable ou formatada em LaTeX.
% - Plotar o gráfico de convergência.
% - TESTE DA SOLUÇÃO: Realizar a avaliação de f(c_final) e atestar a veracidade numérica. Se houver falha, levantar a hipótese do porquê (ex: multiplicidade da raiz par, quebra do teorema de Bolzano, etc).

\subsubsection{Problema Teste 3: Análise de Taxas de Juros em Planos de Financiamento}
\label{subsubsec:problema_teste_3_bisseccao}

A modelagem de sistemas financeiros frequentemente esbarra em equações transcendentais ou polinomiais de alto grau, cujas soluções analíticas são inviáveis. O Problema 3.8 exige a determinação da taxa de juros implícita ($J$) em duas modalidades de financiamento, visando identificar a opção mais vantajosa para o consumidor. O valor financiado ($VF$) corresponde ao preço à vista deduzido da entrada, resultando em $VF = 162.00 - 22.00 = 140.00$.

A priori, a relação matemática que rege o valor presente de uma série de pagamentos uniformes é expressa por:
\begin{equation}
    \frac{1 - (1+J)^{-P}}{J} = \frac{VF}{PM}
    \label{eq:juros_original}
\end{equation}
onde $P$ é o número de parcelas e $PM$ o valor da prestação. No entanto, o isolamento algébrico de $J$ nesta formulação é impossível. Adicionalmente, a Equação \ref{eq:juros_original} apresenta uma singularidade numérica natural quando $J \to 0$. Para contornar essa limitação e condicionar o problema aos métodos de busca de raízes, aplica-se a substituição de variáveis $x = 1+J$ e a constante de proporcionalidade $k = \frac{VF}{PM}$. Dessa forma, o modelo é reescrito na forma polinomial:
\begin{equation}
    f(x) = kx^{P+1} - (k+1)x^P + 1 = 0
    \label{eq:juros_polinomial}
\end{equation}

Para o \textbf{Plano A}, os parâmetros são $P = 9$ e $PM = 26.50$, o que fornece $k_A = \frac{140}{26.5} \approx 5.2830$. A função objetivo torna-se um polinômio de grau 10. Observa-se que a Equação \ref{eq:juros_polinomial} possui uma raiz trivial em $x=1$ (o que implicaria $J=0$, fisicamente inconsistente para um financiamento). Assim, para forçar a convergência do algoritmo para a raiz de interesse financeiro ($J > 0 \implies x > 1$), o intervalo inicial de busca deve ser estritamente superior à unidade.

Os dados processados pelo algoritmo basearam-se no intervalo inicial $[1.05, 1.30]$, que equivale a buscar uma taxa de juros entre $5\%$ e $30\%$ ao período, com uma tolerância rigorosa de $\epsilon < 10^{-6}$. A Tabela \ref{tab:bisseccao_juros} sumariza a evolução da convergência para o Plano A.

\begin{table}[htpb]
\centering
\caption{Histórico iterativo do Método da Bissecção para determinação da taxa do Plano A ($f_1(x)$).}
\label{tab:bisseccao_juros}
\small
\begin{tabular}{@{}cccccc@{}}
\toprule
\textbf{Iteração ($k$)} & $\mathbf{a_k}$ & $\mathbf{b_k}$ & $\mathbf{c_k}$ & $\mathbf{f_1(c_k)}$ & \textbf{Erro} ($\epsilon_k$) \\ \midrule
0  & 1.050000 & 1.300000 & 1.175000 & \phantom{-}0.67780045 & 0.10638297 \\
1  & 1.050000 & 1.175000 & 1.112500 & -0.05891906           & 0.05617977 \\
2  & 1.112500 & 1.175000 & 1.143750 & \phantom{-}0.19422409 & 0.02732240 \\
3  & 1.112500 & 1.143750 & 1.128125 & \phantom{-}0.04375379 & 0.01385041 \\
4  & 1.112500 & 1.128125 & 1.120312 & -0.01301356           & 0.00697350 \\
$\vdots$ & $\vdots$ & $\vdots$ & $\vdots$ & $\vdots$          & $\vdots$   \\
14 & 1.122235 & 1.122250 & 1.122242 & -0.00003841           & 0.00000679 \\
15 & 1.122242 & 1.122250 & 1.122246 & -0.00001211           & 0.00000339 \\
16 & 1.122246 & 1.122250 & 1.122248 & \phantom{-}0.00000104 & 0.00000169 \\
17 & 1.122246 & 1.122248 & 1.122247 & -0.00000553           & 0.00000084 \\ \bottomrule
\end{tabular}
\end{table}

A partir da análise dos dados computacionais, verifica-se que logo na primeira iteração ($k=1$), a detecção de alteração de sinal ($f(a_1) \cdot f(c_1) < 0$) confinou a raiz no subintervalo $[1.1125, 1.1750]$, descartando a parcela superior do domínio. Por mais que a função polinomial possua alto grau — o que frequentemente induz gradientes acentuados e compromete a estabilidade de métodos preditivos como Newton-Raphson —, o Método da Bissecção não demonstrou qualquer sensibilidade à inclinação da curva, atestando sua inquestionável robustez.

Durante as iterações subsequentes, nota-se o declínio assintótico rigoroso do erro. Na 17ª iteração, o critério de parada foi satisfeito, uma vez que o erro relativo computado $\epsilon_{17} \approx 8.49 \times 10^{-7}$ ultrapassou o limiar de tolerância imposto. O valor convergido de $c_{17} \approx 1.122247$ permite extrair diretamente a taxa de juros do Plano A:
\begin{equation}
    J_A = x_A - 1 \approx 0.122247 \quad \text{ou} \quad 12.22\% \text{ ao período.}
\end{equation}

Observa-se ainda que a avaliação do resíduo em $c_{17}$ retornou $f_1(1.122247) \approx -5.53 \times 10^{-6}$, comprovando numericamente que o ponto determinado é, de fato, a raiz do polinômio de financiamento.

\begin{figure}[H]
    \centering
    \begin{tikzpicture}
        \begin{axis}[
            width=14cm, height=7cm,
            title={\bfseries Evolução do Erro e do Resíduo para o Plano A},
            xlabel={Iteração (\(k\))},
            ylabel={Métricas (Escala Logarítmica)},
            ymode=log,
            grid=both,
            grid style={line width=.1pt, draw=gray!20},
            major grid style={line width=.2pt,draw=gray!50},
            legend pos=south west,
            tick label style={font=\footnotesize},
            label style={font=\small},
        ]
        
        \addplot[
            color=cyan, mark=triangle, thick 
        ] table [col sep=comma, x=iter, y=erro] {../dados/saida/questao_3_8/questao_3_8_hist_f1_bisseccao.csv};
        \addlegendentry{Erro \(\epsilon_k\)}

        \addplot[
            color=red, mark=*, mark options={fill=white}, thick 
        ] table [col sep=comma, x=iter, y expr=abs(\thisrow{f(c)})] {../dados/saida/questao_3_8/questao_3_8_hist_f1_bisseccao.csv};
        \addlegendentry{\(|f_1(c_k)|\)}

        \end{axis}
    \end{tikzpicture}
    \caption{Convergência do erro na determinação da taxa (Plano A).}
    \label{fig:convergencia_bisseccao_juros}
\end{figure}

Torna-se evidente que, ao replicar o mesmo arcabouço computacional para o \textbf{Plano B} ($P=12$, $k_B \approx 6.5116$), o método revelará a taxa $J_B$. A tomada de decisão final pelo consumidor deverá ser pautada na comparação quantitativa das soluções extraídas, optando-se pelo plano cuja raiz computada ($x-1$) possua a menor magnitude. A formulação numérica provou-se altamente aderente ao problema proposto, mitigando as singularidades do modelo matemático original por meio da estratégia de substituição polinomial acoplada ao isolamento conservador da Bissecção.

% Para inclusão futura: Repetir a estrutura adotada acima, analisando diferentes condições de contorno ou comportamentos oscilatórios da função que desafiem o Método da Bissecção.

\subsection{Dificuldades Enfrentadas}
\label{subsec:dificuldades_bisseccao}

Durante o ciclo de desenvolvimento, duas barreiras fundamentais requisitaram análise e readequação de código:

A primeira dificuldade decorreu de uma limitação algorítmica associada à manipulação em linguagem Python envolvendo o cálculo de erro relativo. O equacionamento inicial do erro demandava a divisão pela soma dos limites do intervalo $(a_k + b_k)$. Em problemas cujo domínio submetido contemplasse raízes simétricas (como, por exemplo, funções ímpares centradas na origem onde $a_k = -b_k$), o denominador se tornava nulo, causando o encerramento abrupto da aplicação mediante a exceção \texttt{ZeroDivisionError}. A transição adaptativa do erro relativo para o erro absoluto nesses pontos críticos exigiu criteriosa instrumentação de \textit{if-statements}, solucionando a limitação do ambiente sem desrespeitar o rigor numérico.

A segunda dificuldade envolveu a aderência aos padrões de projeto de software estruturado. A intenção de isolar a lógica matemática da escrita de arquivos resultou na necessidade de encapsular o estado iterativo. Em iterações profundas, manter um vetor histórico de dicionários em memória, para apenas no final transferi-los para a rotina de escrita CSV (camada de controle), demonstrou inicialmente um ligeiro sobrepeso na alocação de variáveis. A superação exigiu refinamento na forma como o Python instancia e mapeia essas estruturas temporárias. Ainda que houvesse maior complexidade inicial no planejamento desta abstração, as hipóteses projetuais confirmaram-se válidas, viabilizando o uso de dados de saída padronizados independentemente do método matemático em execução.



